\section{我做错的}

上面那节讲的多数是"发现问题然后改对"，这一节讲三件我判断错了、并且错得有一段时间没意识到的事。

\subsection{语音：连续三次判断错方向}

玩家反馈语音有两个问题：播报完有一声爆鸣，而且第一次开启是普通话、后面变成粤语。

我的第一次判断是"系统里有英文语音在读中文"。这个判断错了。用 \code{say -v '?'} 把系统语音列出来才看到，Eddy、Flo、Grandma 这些名字确实各有真正的中文版本，示例文本就是"你好！我叫 Eddy"，不是英文语音冒充中文。

第二次我改成语音选择逻辑，优先精确匹配 \code{zh-CN}。也错了，实测两个中文语音都还是粤语。

第三次我把爆鸣归因于 \code{speechSynthesis.cancel()} 之后没有等待，加了 300 毫秒的间隔。还是不行。到这一步我才意识到前两次都是在没有证据的情况下改代码——我看到一个现象就猜一个原因，改完让使用者去试，试完再猜下一个。

真正定位问题是从使用者提供的一条信息开始的："一开始是对的，取消一次再开启就又是粤语加爆鸣。"中间唯一发生的事就是 \code{cancel()}。也就是说这个操作会持久破坏后续播报的声音选择，只有刷新页面能恢复；而 \code{cancel()} 在代码里有九个调用点，包括每次出招，所以实际对战中每回合都会触发一次。这解释了"只有第一次是对的"。

改法是把 \code{cancel()} 全部删掉，代价是已经开口的那一句不能中途掐断。但改到这一步，验证的成本已经明显高于这个功能的收益了——音频我没有办法自己验证，无头浏览器没有声卡，每一次判断都要占用使用者的时间。所以我最后的选择是\textbf{整体停用语音}，代码留在开关后面，文档里如实写明停用原因，不再对外宣称具备这个能力。

这件事的教训不是"要小心"，是\textbf{在没有独立验证手段的情况下，不要让使用者反复替我做实验}。第一次失败之后就该做的是把变量拆开做成一个可以一次问清的东西，而不是连续猜。

\subsection{一段 CSS 从来没有被写进文件}

营地页的配招卡片排版很乱，名称、数值和说明挤在同一行。我看着截图以为是设计问题，准备重做布局。

实际原因是那段 CSS 根本没写进样式表。我当时把追加 CSS 的命令用 \code{\&\&} 串在了 \code{node --check} 后面，而那一刻 \code{app.js} 正好因为我刚打错一个括号而报错，于是 \code{cat >> style.css} 被跳过了。后面那句 \code{echo} 照样打印成功，我就以为写进去了。

后果是技能卡退回成 \code{inline-block}，所有文字连成一行。我按"设计问题"去改，改了也没用。

修法很简单，但过程值得记：\textbf{用命令的退出码去判断上一步是否成功，而不是用它的输出}。后面我改成逐个 \code{grep} 确认每个选择器真的落盘了。

\subsection{两次重启后端，两次把密钥弄丢}

密钥按设计只保存在后端进程的内存里，不落盘。我为了加载服务端改动重启了两次，两次都让使用者重新录入。

第一次还可以说是不知道，第二次是明知道的——我甚至提前写了提醒。真正的问题是我没有把"重启会清掉密钥"当成一个需要设计的东西，而是当成一个每次都要打招呼的麻烦。

后来加的改动让这件事好一些：客户端遇到 403 会自动重建会话并重试一次，所以重启之后使用者不会再看到一个莫名其妙的"模型不可用"；错误信息也不再一律显示"暂不可用"，而是把真实原因（会话失效、鉴权失败、超时）带出来。这些都是补丁，根因是这个项目从头到尾没有一个稳定的运行方式。

\subsection{一个用文件名骗过我的数据文件}

有一次我看到 \code{reports/live-model-eval.json} 存在，就以为真实模型评测跑完了。实际上那份文件只有两行，是更早一次局部验证留下的；完整的那次运行还没有做完。

这件事提醒我，交付物里文件名和内容要能对应上，中间产物不要用最终产物一样的名字。后来我在整理清单时对每一条都去看实际数据，而不是看文件名。
